% Section: Methods (drafted 2026-10-30). Every claim cites its source in a comment; every number is \pv{...}.
\section{Methods}
\label{sec:methods}

\subsection{Model and training infrastructure}

% configs/base_4bit.yaml (pinned snapshot, D78); stage_r4.yaml lora; Unsloth, TRL; D70 (Phase 12 trains on the
% 4-bit base); D67, D68 (adapters are never merged for evaluation); FINDINGS #15 (the bf16 merge erased most of the
% RL update; the superseded merged results stay labelled, #13, #14).
All models are \pv{model.base}~\citep{yang2024qwen25math}, loaded with \pv{model.bits}-bit NF4 quantization~\citep{dettmers2023qlora} from a pinned
snapshot and trained with LoRA adapters~\citep{hu2022lora} of rank \pv{lora.r} ($\alpha = \pv{lora.alpha}$) on all
\pv{lora.modules} attention and MLP projections, using Unsloth~\citep{unsloth} and TRL~\citep{vonwerra2020trl}.
Adapters are evaluated as trained: unmerged, on the same quantized base.
An early evaluation path merged reinforcement-learning adapters into a bf16 copy of the model, which erased most of
the update; those results are superseded and kept only as labelled records.
% D137 (the cross-model pilot).
A cross-model pilot repeats the D-pos-v2 and B-v2 recipes of Section~\ref{sec:dpos} on
\pv{model.llama}~\citep{grattafiori2024llama3}.

\subsection{Supervised warm-up}
\label{sec:sft}

% D35 (data and masking), FINDINGS #4 (format and length fixed; the A/B), configs/sft.yaml.
Reinforcement learning starts from a supervised warm-up on \pv{sft.n} generator solutions written in the output
format: the needed SI values, one line per chain step with the substituted numbers, the \texttt{<steps>} block and
the answer.
The $\sigma$ in these targets is a fixed \pv{sft.sigmarel} of the value, and its tokens are masked from the loss, so
the warm-up teaches the format and the computation but never a particular uncertainty.
It runs for \pv{sft.minutes} minutes at learning rate \pv{sft.lr}.

\subsection{Reinforcement learning with a proper scoring rule}
\label{sec:grpo}

% D34, D36 (the GRPO design); configs/grpo_split_credit.yaml (identical to grpo_overnight2.yaml except split credit,
% tested); DeepSeekMath (GRPO); D41 (matched-step comparisons).
We train with GRPO~\citep{shao2024deepseekmath}: for each of \pv{grpo.prompts} problems per step the policy samples
\pv{grpo.G} completions at temperature \pv{grpo.temp}, each completion is scored by the verifier
(Section~\ref{sec:verifier}), and the reward is normalized within its group.
There is no KL penalty ($\beta = \pv{grpo.beta}$); the learning rate is \pv{grpo.lr} and completions are capped at
\pv{grpo.maxlen} tokens.
% CLAUDE.md component 5 (curriculum); grpo_split_credit.yaml curriculum.
Problems are drawn by a curriculum over (family, depth) cells that tracks each cell's pass rate $p$ with an
exponential moving average (rate \pv{grpo.ema}) and samples cells in proportion to $p(1-p) + \pv{grpo.floor}$, the
weight that favours groups with reward variance.
Runs are compared at a matched step (\pv{grpo.step}).

\paragraph{Split credit.}
% Phase 9.5 step 3 (CLAUDE.md); D45, D46 (the trainer); grpo_split_credit.yaml split_credit; D62 (seeds 1, 2);
% D71 / R2-6 (uniform-credit seeds); FINDINGS #6 (the sigma = 0 collapse under sequence-level credit), #16, #18.
With sequence-level credit every token of a completion receives the same advantage, so the tokens of $\sigma$ are
rewarded for the answer's correctness rather than for their own calibration.
Split credit decomposes the verifier's reward $r$ into two parts that sum to it,
\begin{equation}
  r_{\text{answer}} = \crpslog(\mu, \sigma_{\text{ref}}) + w\,r_{\text{steps}}, \qquad
  r_{\sigma} = \crpslog(\mu, \sigma) - \crpslog(\mu, \sigma_{\text{ref}}),
  \label{eq:split}
\end{equation}
with a fixed reference $\sigma_{\text{ref}}$ of \pv{split.ref} of $\mu$ and the step weight $w$ of
Section~\ref{sec:verifier}.
Each part is normalized within its group.
The tokens of the $\sigma$ number, located by parsing, receive the advantage of $r_{\sigma}$; every other token
receives that of $r_{\text{answer}}$.
A failed completion keeps its whole reward in $r_{\text{answer}}$.
The verifier, the logged reward and every hyperparameter are those of the uniform-credit run; split credit and uniform
credit are each trained with \pv{split.seeds} seeds.

\paragraph{$\sigma$-only reinforcement learning (E1).}
% D52 (E1 pre-registration), D94 / D101 (E1c, E1d; E1d's learning rate), R5-7 (E1d is exploratory);
% configs/grpo_e1a.yaml, grpo_e1d.yaml; FINDINGS #14, #25.
To isolate credit assignment completely, E1 holds the answer fixed.
Each training prefix is a warm-up model's own completion truncated right after the ``$\pm$'' of its answer, and the
policy samples only the remaining $\sigma$ suffix (at most \pv{eone.suffix} tokens), so that within a group only
$\sigma$ varies.
E1 trains for \pv{eone.steps} steps at learning rate \pv{eone.lr}; the E1d pilot repeats it at \pv{eone.lrd}, and is
treated as exploratory.

\subsection{Explicit propagation: D-pos-v2 and its implicit control B-v2}
\label{sec:dpos}

% D112, D113 sec. 1 (R4-1 data); results/phase12_sft_data_r4/manifest.json; configs/sft_r4.yaml.
% Motivation: stage B (SFT on sigma*: reads the stated percentages, FINDINGS #27), stage C (sigma-only RL with the
% predictive reward: copies the smallest percentage, #28) and stage D (an explicit line without value-dependent
% derivations: exponent templates, #29) did not propagate.
Supervising $\sigstar$ directly, rewarding $\sigma$ alone, and training an explicit propagation line on data without
value-dependent derivatives all failed to make the stated $\sigma$ propagate noise.
The data of this section contain the behaviour.
They reuse \pv{rfour.n} training problems with mixed noise: each given's level is drawn independently from
\pv{gen.mixedlevels}, and \pv{rfour.mixing} problems mix two or more noisy levels.
Both arms see the same problems.

% D113 sec. 1 (the explicit arm's format and masks); sft_r4.yaml arms.
\paragraph{D-pos-v2 (explicit).}
After the solution, the target completion writes one line per noisy given $x_i$:
the given's total log-sensitivity $S_i = \partial\ln f/\partial\ln x_i$ as a symbolic expression derived by the chain
rule over the steps, the same expression with the solution's values substituted, and its value to
\pv{rfour.sfmt} significant figures.
A power-law term is stated as its exponent.
A propagation line then combines the terms with the stated levels $s_i$,
\begin{equation}
  \sigma_{\text{rel}} = \sqrt{\textstyle\sum_i (S_i\, s_i)^2} = X,
  \label{eq:line}
\end{equation}
and the answer states $\sigma = X\,|\mu|$.
Every token is trained, including the sensitivities, the line and $\sigma$.
% D113 sec. 3 (the three sigma values; (c_S) descriptive).
We read three values of $\sigma$ from each D-pos-v2 completion:
(a) the answer's relative $\sigma$;
(b) the line's stated $X$;
(c) the exact value of Equation~\eqref{eq:line} computed from the line's own written terms.
Comparing them locates a failure in the propagation itself, the arithmetic of the line, or the copy into the answer.

% D113 sec. 1 (the supervised arm), D104 (stage B's arm).
\paragraph{B-v2 (implicit).}
The same problems and solutions with no sensitivity lines and no propagation line; $\sigstar$ itself is written after
``$\pm$'' and trained.
B-v2 can learn propagation only implicitly.

% D113 sec. 2, D118 sec. 1 (recipe; seeds change the LoRA initialization, data order and trainer seed; data fixed);
% stage_r4.yaml; D114 (token budget); D116 (eval batch).
\paragraph{Training and evaluation.}
Each arm trains a fresh adapter on the pinned base for \pv{rfour.epochs} epochs (\pv{rfour.steps} steps) at
learning rate \pv{rfour.lr} and effective batch \pv{rfour.batch}, with \pv{rfour.seeds} seeds per arm.
The seed changes the adapter's initialization and the data order, not the data.
Evaluation is greedy, with up to \pv{rfour.tokens} new tokens, enough for the longest derivations.

\subsection{Reference uncertainties}
\label{sec:comparators}

% D95 (the noise heuristics; RSS over the givens' stated levels is the primary); physcal/eval/noise_heuristics.py.
Every stated $\sigma$ is compared with three references that need no model.
$\sigstar$ is the generator's own value (Equation~\eqref{eq:sigstar}).
RSS is the root sum of squares of the stated levels of the givens, $\sqrt{\sum_i s_i^2}$: correct for a product of
first powers, and blind to the values.
% D73 (X1 pre-registration), D99 (exact refit, reproduction gate before use), FINDINGS #17, #23.
X1 is a shallow text model of $\sigstar$: gradient-boosted trees on the problem text's tokens, the stated levels and
the displayed magnitudes, fitted on generator problems disjoint from every evaluation set.
It is refit exactly and used only if it reproduces its published results.
% D147 sec. 10 (the clamped comparator, J set only).
On the J set we add a fourth, \emph{clamped} reference: the first-order $\sigma$ with every sensitivity clipped into
its chain's training range, which is what a model that propagates exactly but only within its training range would
state.

\subsection{Executing instead of predicting: MC-A and MC-B}
\label{sec:mc}

% D112 (R4-4), D113 sec. 7 (MC-A, inference only); configs/mc_a.yaml; FINDINGS #31.
MC-A estimates $\sigma$ by re-asking the model.
For each problem it builds \pv{mca.k} perturbed copies in which every noisy given is redrawn from the generator's own
predictive distribution, $x = x_{\text{meas}}/(1+s z)$ with $z \sim \mathcal{N}(0,1)$, and shown in its displayed
unit; B-v2 answers each copy.
Its $\sigma$ is the spread of $\ln\mu$ over the copies: the standard deviation (raw), or \pv{mca.mad} times the median
absolute deviation over the copies whose solution steps match the unperturbed one (robust).
The same perturbations pushed through the exact chain give an oracle that targets $\sigstar$.

% D119 (MC-B pre-registration), D122 (deviations while building; the token budget); configs/mc_b.yaml;
% FINDINGS #33, #37; PAL and Program of Thoughts limit its novelty.
MC-B asks the model for the computation instead of the answer, in the spirit of program-aided
reasoning~\citep{gao2023pal,chen2023pot}.
The base model, without an adapter, writes a Python function \texttt{solve} whose parameters are the givens' own
symbols; \pv{mcb.nexamples} worked examples from a disjoint seed range at depth \pv{mcb.exdepth} show the format, and
generation is greedy with up to \pv{mcb.tokens} new tokens.
A sandbox checks the program's syntax tree (only the \texttt{math} module, no dynamic code), calls it on the
displayed values and on MC-A's perturbations, and takes the spread of the logarithms of its outputs as $\sigma$.
A failing program is classified as an interface failure (no program, wrong signature, rejected), a runtime failure,
or a semantic failure (it runs but computes the wrong chain).
